\begin{longtable}{
    |>{\RaggedRight\arraybackslash}p{0.319\textwidth}
    |p{0.637\textwidth}|
    } % 0.956
    \caption{Определения цикла в учебной литературе}
    \label{tab:cycle-definitions} \\
    \LongTableHeader{
        \hline
        \rowcolor{black}
        \textcolor{white}{\textbf{Источник}} &
        \textcolor{white}{\textbf{Определение цикла}} \\
        \hline
    }

    \longcaptioneachpage{Источник: составлено автором}

    Mankiw N. G. (2010). Macroeconomics, 7th ed. Worth Publishers\cite{Mankiw-2010}
    &
    Business cycles – \textquotedbl{}short-run fluctuations in economic activity around its long-run trend\textquotedbl{}; понимаются как отклонения реального ВВП, занятости и других показателей от нормального уровня \\ \hline

    Blanchard O. (2016). Macroeconomics, 7th ed. Pearson\cite{Blanchard-2016}
    &
    Экономические флуктуации (business cycle fluctuations) – отклонения выпуска, занятости, потребления и инвестиций от долгосрочного тренда \\ \hline

    Dornbusch R., Fischer S., Startz R. (2011). Macroeconomics, 11th ed. McGraw-Hill\cite{Dornbusch-2011}
    &
    Business cycles are the recurring fluctuations that occur in real GDP over time \\ \hline

    Samuelson P. A., Nordhaus W. D. (2010). Economics, 19th ed. McGraw-Hill\cite{Samuelson-2010}
    &
    Деловой цикл – периодическое, но нерегулярное колебание деловой активности: за расширением экономики следует спад, за спадом – оживление; ключевые индикаторы – ВВП, занятость, доходы \\ \hline

    Чепурин М. Н., Киселёва Е. А. (2006). Курс экономической теории. 5-е изд. Киров: АСА\cite{Чепурин-2006}
    &
    Экономический цикл – это повторяющийся на протяжении ряда лет колебания различных показателей экономической активности: темп роста ВНП, общий объём продаж, общий уровень цен, уровень безработицы, загрузка производственных мощностей, величина инвестиций»; основным показателем цикличности служат колебания темпа роста ВНП \\ \hline

    Тарасевич Л. С., Гребенников П. И., Леусский А. И. (2010). Макроэкономика: учебник\cite{Тарасевич-2010}
    &
    Экономический (деловой) цикл – периодически повторяющиеся спады и подъёмы в экономике, колебания деловой активности; колебания нерегулярны и непредсказуемы, поэтому термин \textquotedbl{}цикл\textquotedbl{} условен \\ \hline

    Агапова Т. А., Серегина С. Ф. (2004). Макроэкономика. 6-е изд. М.: Дело и сервис\cite{Агапова-2004}
    &
    Экономический (деловой) цикл – подъёмы и спады уровней экономической (деловой) активности в течение нескольких лет; промежуток времени между двумя одинаковыми состояниями экономической активности \\ \hline

    Аносова А. В., Ким И. А., Серегина С. Ф. и др. (2013). Макроэкономика: учебник для бакалавров. 2-е изд. М.: Юрайт\cite{Аносова-2013}
    &
    Экономический цикл описывается как последовательность фаз подъёма, пика, спада и дна; подчёркивается отклонение фактического ВВП от потенциального и нерегулярность колебаний \\ \hline
\end{longtable}